\documentclass[11pt]{article}

% ── Packages ──────────────────────────────────────────────────────────────────
\usepackage[margin=1in]{geometry}
\usepackage{amsmath,amssymb}
\usepackage{booktabs}
\usepackage{graphicx}
\usepackage[hidelinks]{hyperref}
\usepackage{natbib}
\usepackage{microtype}
\usepackage{caption}
\usepackage{enumitem}

\bibliographystyle{plainnat}

\title{\textbf{Quote Placement Dominates Microstructure Prediction\\
in Retail-Latency Market Making}\\[0.4em]
\large A Negative Result from 70 Days of Cryptocurrency Limit Order Book Data}

\author{Jude Kriel Ramcharitar\\
\small Independent Researcher\\
\small \texttt{j.k.ramcharitar@gmail.com}}

\date{August 2026 \quad$\cdot$\quad Version 2.0 --- Working Paper}

\begin{document}
\maketitle

% ══════════════════════════════════════════════════════════════════════════════
\begin{abstract}
\noindent
High-frequency market makers profit by acting on price signals faster than
competitors. Retail and non-institutional algorithmic traders, operating under
execution delays of 200\,ms to 2\,s, are effectively excluded from market making:
by the time a quote reaches the exchange, the market has moved. A natural
conjecture, and the premise of an earlier version of this work, is that
machine learning models trained on limit order book (LOB) microstructure can
substitute for speed by anticipating market state at order arrival.

We test that conjecture and reject it at retail-accessible data frequency. Using
269{,}330 LOB snapshots collected continuously over 70 days from Kraken across
BTC/USD, ETH/USD, and SOL/USD, we train a survival-analysis fill probability
model on 109 microstructure features and compare it against a deliberately
trivial benchmark: a 32-cell lookup table containing only the mean historical
fill rate for each combination of quote offset, latency window, and side. On a
matched test split, the neural model attains AUC 0.6375 and the lookup table
0.6372, a difference of $+0.0003$. Per-asset replication reproduces the null
on all three instruments, On the smallest asset, ETH/USD, the neural model \emph{underperforms} the
table by 0.0988.

We report two further findings. First, a controlled ablation shows the entire
measurable predictive content of the model is carried by three scalar quantities
describing the quote itself; zeroing the 109-dimensional LOB window and
retaining only those scalars costs 0.085 AUC, while the reverse configuration
collapses to chance ($0.5034 \pm 0.031$ across five seeds). Second, we document
a conditioning defect that invalidated our own earlier results, in which quote
offset was omitted from the model input, rendering four training rows with fill
rates spanning 0.782 to 0.358 numerically identical at the input layer. We
report this explicitly, and treat the corrected pipeline as a reproduction
control throughout.

The practical implication is direct. At roughly 90 LOB updates per hour, the
rate available from a standard public exchange WebSocket feed, fill probability
is a function of how far from mid a trader quotes, and microstructure signals
add nothing measurable on top. This bounds where the necessary signal can be
found rather than refuting latency compensation as a principle, and it indicates
that the binding constraint for retail participants is data frequency rather
than model capacity or sample size.

\vspace{0.8em}
\noindent\textbf{Keywords:} market making, market microstructure, limit order
book, fill probability, execution latency, retail trading, negative results,
survival analysis

\vspace{0.3em}
\noindent\textbf{JEL:} G12, G14, C45, C61, G23
\end{abstract}

\vspace{1em}

% ══════════════════════════════════════════════════════════════════════════════
\section{Introduction}

In competitive markets, speed is capital. Firms dominating modern equity and
cryptocurrency liquidity provision have invested heavily in compressing
execution latency to microseconds. Microseconds do not matter to the real
economy; being first to respond to a price signal is profitable, repeatedly, at
scale. This raises a precise question: is speed
itself the valuable resource, or is it a proxy for the ability to act on
information before it becomes stale?

If the latter, speed might not be strictly necessary. A participant who cannot
act in microseconds could still act correctly, provided they can predict what
the market will look like when their order arrives. Version~1 of this work
proposed exactly that: \emph{latency-compensating market making}, in which deep
learning models trained on LOB microstructure substitute predictive accuracy for
execution speed.

This version reports that the conjecture does not survive contact with the data
we are able to collect, and characterises precisely why.

\subsection{What changed}

Version~1 presented a theoretical framework alongside results on synthetically
generated LOB data. Those synthetic results were encouraging. They were also,
we now recognise, uninformative: a geometric Brownian motion price process with
a hand-specified fill rule contains whatever structure the specification puts
there, and the fill probability model was in effect recovering its own
generator.

On real data the picture inverted. Our first real-data results were worse than
chance, which prompted a diagnostic that located a defect in the model's input
pipeline (Section~\ref{sec:defect}). Correcting it produced a genuine
improvement. But a like-for-like comparison against a trivial baseline (the comparison we
should have run first) showed that the corrected model does not outperform a
lookup table on three variables.

We report the sequence in full, including our own error, because the error is
instructive: a model can appear to improve with more data, respond to ablation
as though its features matter, and still carry no information beyond a trivial
benchmark that was never run.

\subsection{Contributions}

\begin{enumerate}[itemsep=2pt,topsep=4pt]
\item We formalise the latency-compensating market making problem as an
extension of \citet{avellaneda2008} incorporating execution delay
(Section~\ref{sec:formal}). This framework is unchanged from Version~1 and we
believe it remains correct; our negative result concerns its empirical
instantiation, not its logic.

\item We release 269{,}330 timestamped 10-level LOB snapshots across three
cryptocurrency pairs, collected continuously over 70 days, together with the
complete collection and feature pipeline (Section~\ref{sec:data}).

\item We show that a survival-analysis fill probability model over 109
microstructure features does not outperform a 32-cell lookup table on quote
offset, latency, and side, either pooled ($+0.0003$ AUC) or per-asset
(Section~\ref{sec:results}).

\item We isolate the source of predictive content via a three-arm ablation,
showing that quote geometry accounts for essentially all of it
(Section~\ref{sec:ablation}).

\item We document a model conditioning defect that invalidated our own prior
results, and use its deliberate reproduction as an experimental control
(Section~\ref{sec:defect}).
\end{enumerate}

\subsection{What this result does and does not establish}
\label{sec:scope}

Because negative results are easy to over-read, we state the scope precisely.

\textbf{Established.} At an observed update frequency of approximately 90 LOB
snapshots per hour per asset (the rate delivered by Kraken's public WebSocket
book channel at depth 10), a 90{,}754-parameter LSTM over 109
engineered microstructure features provides no measurable improvement in fill
prediction over quote offset, latency window, and side alone, on any of three
cryptocurrency pairs.

\textbf{Not established.} We have \emph{not} shown that LOB microstructure is
uninformative about fill probability in general. A substantial literature
demonstrates otherwise at higher frequencies
\citep{maglaras2022,arroyo2024}. Our data are sampled at a median inter-update
interval of 18.2 seconds (mean 40.1 seconds), which is three to four orders of
magnitude coarser than the tick-level data used in that work. The most probable
explanation for our null is that microstructure signal exists but does not
survive this sampling interval. We regard distinguishing that explanation from
the alternatives as the central open question
(Section~\ref{sec:interpretation}).

\textbf{Not tested.} The reinforcement learning agent proposed in Version~1
consumed fill probability estimates from the defective model and was trained
only on synthetic data. We therefore report no RL results in this version. See
Section~\ref{sec:rl-status}.

% ══════════════════════════════════════════════════════════════════════════════
\section{Related Work}

\subsection{Market making theory and latency}

The modern framework originates with the inventory models of \citet{garman1976}
and \citet{ho1981}, formalised as a stochastic control problem by
\citet{avellaneda2008}, who derived closed-form optimal quotes as functions of
inventory, risk aversion, and order arrival intensity. Extensions incorporate
adverse selection and execution uncertainty \citep{cartea2018}, and
\citet{cartea2023} introduce stochastic delay into optimal execution.

Nearly all analytical market making models assume negligible execution latency.
\citet{lehalle2022} show this assumption is load-bearing rather than
cosmetic: the value of predictive signals about future order flow is
systematically eroded by latency, since a trader who can predict adverse flow
but cannot act quickly enough derives diminishing benefit. Our result is consistent with theirs and sharpens it. In our setting the
erosion is complete at the data frequency available to us.

\subsection{The latency arms race}

\citet{moallemi2013} provide the first formal model of latency cost.
\citet{budish2015} characterise continuous-time limit order books as generating
structurally wasteful speed investment. \citet{aquilina2022} quantify this
empirically, finding latency arbitrage races occur roughly once per minute per
symbol in FTSE 100 stocks with modal durations of 5--10 microseconds, and
estimating that eliminating them would reduce effective spreads by
approximately 17\%.

Proposed remedies (the IEX speed bump of \citealp{hu2019}, frequent batch
auctions in \citealp{budish2015}, micro-burst fees in \citealp{brolley2023})
operate at the exchange level and require regulatory or venue adoption. None offers
recourse to an individual trader on existing infrastructure, which was the
motivation for pursuing a participant-level technical approach.

\subsection{Fill probability estimation}

\citet{lo2002} established econometric foundations for execution probability
modelling. \citet{maglaras2022} introduced a deep learning approach treating
execution as a time-to-fill distribution. \citet{arroyo2024} advanced this with
a convolutional-transformer architecture using survival analysis, which handles
the right-censoring inherent in limit order data where many orders are cancelled
before executing.

Our fill probability model follows \citet{arroyo2024} in its survival
formulation. The critical difference in our setting is data frequency: that work
uses tick-level message data, whereas we use snapshot data at a median 18.2
second interval. We return to this in Section~\ref{sec:interpretation}, as we
believe it is the primary explanation for our divergent result.

\subsection{Reinforcement learning for market making}

\citet{spooner2018} demonstrated RL agents learning competitive market making
policies in simulation; \citet{guo2023} added continuous action spaces and
attention-based LOB features; \citet{briola2021} established viability in
high-frequency cryptocurrency settings. The most directly relevant work is
RELAVER \citep{jiang2025}, which explicitly models random execution delays of
30--100\,ms with a batch matching mechanism.

Version~1 of this paper positioned itself as extending RELAVER to a retail
latency regime an order of magnitude larger. We no longer make that claim, as
the predictive component such an extension would require does not function at
our data frequency.

% ══════════════════════════════════════════════════════════════════════════════
\section{Problem Formalisation}
\label{sec:formal}

We retain the formalisation from Version~1, as the negative result concerns its
empirical instantiation rather than its structure.

\subsection{Standard setup}

A market maker operates over horizon $[0,T]$. The mid-price follows
$S_t = S_0 + \sigma W_t$ with $\sigma > 0$ and $W_t$ standard Brownian motion.
Posting a bid at distance $\delta^b$ and ask at $\delta^a$ from mid, fills
arrive as Poisson processes with intensities
\begin{equation}
\lambda^b(\delta^b) = A e^{-k\delta^b}, \qquad
\lambda^a(\delta^a) = A e^{-k\delta^a},
\end{equation}
where $A, k > 0$ govern the decay of arrival rate with distance from mid.
Inventory $q_t$ evolves with fills and wealth $X_t$ accumulates spread capture
net of inventory risk. The objective is
\begin{equation}
\max_{\delta^b,\delta^a} \; \mathbb{E}\!\left[ X_T + q_T S_T - \gamma q_T^2 \right],
\end{equation}
with $\gamma > 0$ the inventory risk aversion parameter.

\subsection{Introducing execution latency}

We depart from the standard framework by introducing execution latency $\ell$,
the delay between quoting decision and order arrival at the exchange. We
consider $\ell \in [200\,\text{ms}, 2000\,\text{ms}]$, characterising traders
without co-location infrastructure.

Under latency $\ell$, a quote posted at time $t$ arrives at $t + \ell$ into an
order book state $\mathcal{F}_{t+\ell}$ that is unobserved at decision time. The
relevant fill probability is
\begin{equation}
P^{\text{fill}}(\delta, \ell, \mathcal{F}_t)
= \mathbb{P}\!\left[\text{order fills} \mid \delta, \mathcal{F}_{t+\ell}\right],
\end{equation}
which depends on a future state. This creates two compounding losses relative to
the zero-latency benchmark: \emph{adverse selection amplification}, where quotes
are picked off when the mid moves against them during transmission, and
\emph{cancellation impotence}, where a trader detecting an adverse signal cannot
cancel before being filled.

\subsection{The latency-compensating problem}

Let $\widehat{P}^{\text{fill}}$ denote an estimator of fill probability and
$\widehat{S}_{t+\ell}$ an estimator of the mid-price at arrival. The
latency-compensating problem is
\begin{equation}
\max_{\delta^b,\delta^a} \;
\mathbb{E}\!\left[ X_T + q_T \widehat{S}_T - \gamma q_T^2
\;\middle|\; \widehat{\mathcal{F}}_{t+\ell} \right]
\end{equation}
subject to the participation constraints
\begin{equation}
\widehat{P}^{\text{fill}}(\delta, \ell, \mathcal{F}_t) \geq \tau_{\text{fill}},
\qquad
\left| \widehat{S}_{t+\ell} - S_t \right| \leq \Delta_{\text{safe}},
\label{eq:abstain}
\end{equation}
with abstention whenever either constraint is violated.

\subsection{Viability conditions}
\label{sec:conditions}

Version~1 stated three conditions for viability. We restate them, since our
result speaks directly to the first.

\begin{description}[itemsep=3pt,topsep=4pt]
\item[Condition 1 --- Predictive accuracy.] The estimators
$\widehat{P}^{\text{fill}}$ and $\widehat{S}_{t+\ell}$ must be accurate enough
that the constraints in \eqref{eq:abstain} select genuinely favourable states
rather than arbitrary ones.

\item[Condition 2 --- Uninformed order flow.] The market must contain
sufficient noise-trader flow to generate fill opportunities that are not
systematically adverse \citep{chen2018}.

\item[Condition 3 --- Abstention discipline.] The abstention rule must be
strictly enforced; a strategy that quotes indiscriminately degrades to the
latency-penalised baseline.
\end{description}

\noindent
Our empirical finding is that Condition~1 fails at retail-accessible data
frequency. The failure is not that no estimator achieves discrimination. It is
that the achievable discrimination is entirely reproduced by a lookup table on
quote geometry, which requires no model and confers no informational
advantage.
Conditions 2 and 3 are untested here.

% ══════════════════════════════════════════════════════════════════════════════
\section{Data}
\label{sec:data}

\subsection{Collection}

LOB data were collected continuously from Kraken via WebSocket API v2,
subscribing to the \texttt{book} channel at depth 10 for BTC/USD, ETH/USD, and
SOL/USD. Each update is timestamped at receipt and stored with the full 10-level
bid and ask ladder in Parquet format, partitioned by asset, date, and hour.

Collection ran on a cloud virtual machine from 15 June to 24 August 2026. The
collector is managed by a systemd unit with \texttt{Restart=always}; over the
final three-week window it was restarted automatically three times following
WebSocket disconnections. An earlier configuration using a detached terminal
session and a bounded reconnect loop failed silently on 20 July and remained
down for fourteen days, costing an estimated 90{,}000 snapshots. We note this
because unattended collection reliability is a material and under-discussed
constraint on independent microstructure research.

\subsection{Sampling characteristics}

Kraken's book channel is event-driven: updates are emitted when the visible book
changes, not on a fixed schedule. The resulting sampling is both sparse and
irregular, and this is central to interpreting our results.

\begin{table}[h]
\centering
\caption{Inter-update intervals, BTC/USD.}
\label{tab:gaps}
\begin{tabular}{lr}
\toprule
Statistic & Value \\
\midrule
Mean gap            & 40.1\,s \\
Median gap          & 18.2\,s \\
Minimum gap         & $<$0.01\,s \\
Maximum gap         & 6339\,s \\
Updates per hour    & $\approx$90 \\
\bottomrule
\end{tabular}
\end{table}

\noindent
The maximum gap corresponds to a collector outage rather than market quiescence.
The mean-median divergence indicates a heavy right tail: most updates cluster,
with long quiet intervals between bursts.

This sampling rate is the single most important feature of our dataset. A
latency window of 200\,ms is roughly one ten-thousandth of the mean interval
between observations. The model is therefore asked to discriminate outcomes at a horizon far shorter
than its own sampling resolution. We return to this in
Section~\ref{sec:interpretation}.

\begin{table}[h]
\centering
\caption{Collected dataset, 15 June -- 24 August 2026.}
\label{tab:data}
\begin{tabular}{lrrr}
\toprule
Asset & Snapshots & Label rows & Overall fill rate \\
\midrule
BTC/USD & 111{,}002 & 3{,}551{,}392 & 0.734 \\
ETH/USD &  34{,}380 & 1{,}099{,}488 & 0.754 \\
SOL/USD & 123{,}948 & 3{,}965{,}664 & 0.534 \\
\midrule
Total   & 269{,}330 & 8{,}616{,}544 & --- \\
\bottomrule
\end{tabular}
\end{table}

\noindent
ETH/USD yields substantially fewer snapshots over an identical collection
window. The visible book on this venue is genuinely less active; this is not a
collection artefact. The asymmetry matters for Section~\ref{sec:perasset}.

\subsection{Label construction}

For each snapshot we simulate limit order submissions at quote offsets
$\delta \in \{1,2,3,5\}$ ticks on each side, and record whether each simulated
order would have filled within each latency window
$\ell \in \{200, 500, 1000, 2000\}$\,ms based on subsequent book evolution. This
yields $4 \times 4 \times 2 = 32$ label rows per snapshot. Orders not filled
within $\ell$ are treated as right-censored at $\ell$.

Table~\ref{tab:marginals} reports marginal fill rates. Both variables show strong, monotone structure, which is what makes the
lookup-table baseline of Section~\ref{sec:baseline} a demanding comparison
rather than a straw one.

\begin{table}[h]
\centering
\caption{Marginal fill rates by quote offset and latency window, SOL/USD.}
\label{tab:marginals}
\begin{tabular}{lr@{\hskip 2.5em}lr}
\toprule
Offset (ticks) & Fill rate & Latency (ms) & Fill rate \\
\midrule
1 & 0.782 & 200  & 0.385 \\
2 & 0.550 & 500  & 0.498 \\
3 & 0.446 & 1000 & 0.585 \\
5 & 0.358 & 2000 & 0.669 \\
\bottomrule
\end{tabular}
\end{table}

\subsection{Features}

From each snapshot we construct 109 features: 40 raw price and volume levels,
normalised prices and volumes, mid-price, spread, spread in basis points, order
book imbalance, depth imbalance, and rolling statistics and lagged differences of
these quantities. Prices are normalised by contemporaneous mid; volumes by a
rolling mean.

Two pipeline parameters required adjustment for sparse data. The rolling
normalisation window was reduced from 1000 to 10 snapshots, and lag features
from $\{1,5,10,20,50\}$ to $\{1,2,3,5\}$; at the original settings, nearly all
rows were eliminated by the subsequent \texttt{dropna}, which we replaced with
zero-fill. These changes are necessary consequences of the sampling rate: a
1000-snapshot rolling window spans roughly eleven hours of wall-clock time at 90
updates per hour, which is not a meaningful local normalisation.

% ══════════════════════════════════════════════════════════════════════════════
\section{Fill Probability Model}
\label{sec:model}

\subsection{Survival formulation}

We model time-to-fill $T \geq 0$ with survival function
$S(t) = \mathbb{P}[T > t]$ and hazard $h(t)$. The fill probability within the
latency window is
\begin{equation}
\widehat{P}^{\text{fill}}(\delta, \ell) = 1 - S(\ell) = \mathbb{P}[T \leq \ell].
\end{equation}
Observations are right-censored: we observe $\widetilde{T} = \min(T, C)$ with
censoring indicator $\Delta$. Models are trained by minimising the negative
log-likelihood under censoring,
\begin{equation}
\mathcal{L} = -\frac{1}{N}\sum_{i=1}^{N}
\left[ \Delta_i \log h(\widetilde{T}_i) + \log S(\widetilde{T}_i) \right].
\end{equation}

\subsection{Architecture}

The model encodes a lookback window of $L = 50$ snapshots with a two-layer LSTM
(hidden dimension 64, dropout 0.2), producing a representation $h$. A
conditioning vector
\begin{equation}
c = \left[ \ell / \ell_{\max}, \; \delta / \delta_{\max}, \; \text{side} \right]
\in \mathbb{R}^3
\end{equation}
is passed through a small MLP and concatenated with $h$ before a monotone
survival decoder with $K$ discrete hazard outputs. Total parameters: 90{,}754.
Training uses Adam at $\eta = 10^{-3}$ with early stopping on validation loss.

The composition of $c$ is the subject of Section~\ref{sec:defect}.

\subsection{Evaluation}
\label{sec:metrics}

We report two ranking metrics that answer different questions, and we report
both throughout.

\textbf{Fill AUC.} Area under the ROC curve for the binary filled/not-filled
task, computed exactly via the Mann-Whitney rank-sum identity with ties credited
0.5. This measures whether the model ranks orders that filled above orders that
did not.

\textbf{Concordance index.} Computed on filled orders only, measuring whether
predicted fill probability ranks consistently with realised fill \emph{time}.

AUC is the metric of primary interest, because the abstention rule in
\eqref{eq:abstain} tests $\widehat{P}^{\text{fill}} \geq \tau_{\text{fill}}$,
a binary fill/no-fill decision that never compares timings. We fixed this
choice on the argument above, prior to observing that AUC would be the
\emph{lower} of the two numbers on our data; we report both regardless.

% ══════════════════════════════════════════════════════════════════════════════
\section{A Conditioning Defect and Its Use as a Control}
\label{sec:defect}

Our first real-data results were substantially worse than chance. Our initial hypothesis was insufficient data. That hypothesis was wrong, and
the diagnostic that displaced it is worth describing, because both the error and
its detection are informative.

\subsection{The defect}

The conditioning vector $c$ originally contained only $\ell$. Quote offset
$\delta$ and side were used to construct labels but never reached the model.

The consequence is structural. For a fixed snapshot and latency, the eight label
rows spanning four offsets and two sides present \emph{identical} model input
while carrying different targets, with marginal fill rates ranging from 0.782 at
1 tick to 0.358 at 5 ticks (Table~\ref{tab:marginals}). No function of the
input can separate them. The loss-minimising response is to predict the
conditional mean and ignore the input entirely, which is what the model did:
validation loss moved by roughly 0.003 over 40 epochs across every run.

This also explains an observation that had misled us. Across successive
collection milestones, measured C-index moved 0.5008, 0.5308, 0.5526 as the
dataset grew from 21K to 51K to 103K snapshots. We read this as convergence
toward a target with increasing data. It was not: all three values lie within
seed variance of chance, and the apparent trend was noise. A subsequent run on 269K snapshots returned 0.4749, worse with more data. That
is the expected behaviour when additional data supplies additional mutually
contradictory input-target pairs.

\subsection{Correction and control}

Adding $\delta$ and side to $c$ raised C-index from 0.4749 to 0.6266 on an
otherwise identical configuration, and validation loss movement increased
roughly eightfold (0.003 to 0.023).

We retain the defective configuration as an experimental arm, denoted
\textsc{lob}, in which the LOB window is supplied but $\delta$ and side are
reset to constants. It is a reproduction control. Any claim that a result derives from LOB features
should come with a demonstration that the corresponding defective configuration
performs at chance. In our ablation it
does, at $0.5034 \pm 0.031$ across five seeds
(Section~\ref{sec:ablation}).

\subsection{Consequence for prior results}

All real-data figures reported in Versions 1.0--1.7 of this paper were produced
by the defective configuration and should be disregarded. This includes the
C-index progression cited above and the projection that approximately 300{,}000
snapshots would be required for convergence, which was extrapolated from noise.

% ══════════════════════════════════════════════════════════════════════════════
\section{Results}
\label{sec:results}

\subsection{The baseline}
\label{sec:baseline}

We compare the model against a lookup table: the mean fill rate for each
$(\delta, \ell, \text{side})$ cell, 32 cells in total, estimated on training
data and applied unchanged to test data. It has no learned parameters, no access
to the order book, and no temporal information.

We emphasise that this is a demanding baseline rather than a straw one. The
marginals in Table~\ref{tab:marginals} show quote offset spanning a 0.42 range
in fill rate and latency spanning 0.28. Any model that fails to beat this table
is failing to add value over two variables the trader chooses directly and
already knows.

\subsection{Pooled result}

\begin{table}[h]
\centering
\caption{LSTM versus lookup table, pooled across three assets, matched test
split ($n = 29{,}983$).}
\label{tab:baseline}
\begin{tabular}{lrr}
\toprule
Predictor & Parameters & AUC \\
\midrule
Lookup table (32 cells) & 32 & 0.6372 \\
LSTM (109 LOB features) & 90{,}754 & 0.6375 \\
\midrule
Difference & & $+0.0003$ \\
\bottomrule
\end{tabular}
\end{table}

\noindent
The difference is three ten-thousandths of an AUC point. For reference, seed
variance alone on this configuration is $\pm 0.034$, two orders of magnitude
larger. The neural model and the lookup table are, on this evidence,
indistinguishable.

\subsection{Per-asset replication}
\label{sec:perasset}

We repeated the comparison independently per asset, training separate models and
fitting separate tables, with three seeds each. Decision thresholds were fixed
before running: a gap exceeding $+0.02$ would indicate LOB signal; within
$\pm 0.02$, none.

\begin{table}[h]
\centering
\caption{Per-asset comparison, AUC, three seeds per asset.}
\label{tab:perasset}
\begin{tabular}{lrrrr}
\toprule
Asset & Snapshots & Lookup & LSTM (mean $\pm$ sd) & Gap \\
\midrule
BTC/USD & 111{,}002 & 0.6396 & $0.6379 \pm 0.006$ & $-0.0017$ \\
ETH/USD &  34{,}380 & 0.6510 & $0.5522 \pm 0.005$ & $-0.0988$ \\
SOL/USD & 123{,}948 & 0.6821 & $0.6739 \pm 0.012$ & $-0.0082$ \\
\bottomrule
\end{tabular}
\end{table}

\noindent
No asset shows LOB signal. Three observations follow.

First, SOL/USD carries the highest lookup AUC in the study, 0.6821. Its fill
rate falls from 0.654 at 1 tick to 0.251 at 5 ticks at 200\,ms latency, a far
steeper gradient than BTC or ETH. We had hypothesised this thinner, more
competitive book would be where microstructure mattered most. Instead the
steepness is fully captured by quote geometry, leaving nothing for the model to
add.

Second, ETH/USD is a genuine failure rather than a null: the model
underperforms the table by 0.099, consistently across seeds (sd 0.005). Its
training runs terminated at epochs 12, 9, and 29 with validation loss moving
only 1.7437 to 1.7316. With 34{,}380 snapshots, roughly a third of the other
assets, the model is undertrained and actively discards information the table
retains. We report this rather than tuning it away, as it bounds the sample size
below which this architecture is counterproductive.

Third, the pooled null is not an artefact of averaging heterogeneous assets. It
holds separately on each.

\subsection{Ablation}
\label{sec:ablation}

To locate the source of predictive content we ran three arms with identical
architecture, data, and schedule, varying only the input:

\begin{description}[itemsep=2pt,topsep=4pt]
\item[\textsc{full}] LOB window plus $[\ell, \delta, \text{side}]$.
\item[\textsc{cond}] LOB window zeroed; conditioning retained.
\item[\textsc{lob}] LOB window retained; $\delta$, side reset to constants (the
defective configuration of Section~\ref{sec:defect}).
\end{description}

\begin{table}[h]
\centering
\caption{Ablation, C-index, five seeds per arm.}
\label{tab:ablation}
\begin{tabular}{lrrrr}
\toprule
Arm & Mean & SD & Min & Max \\
\midrule
\textsc{full} & 0.6101 & 0.034 & 0.5670 & 0.6536 \\
\textsc{cond} & 0.5247 & 0.041 & 0.4853 & 0.5885 \\
\textsc{lob}  & 0.5034 & 0.031 & 0.4678 & 0.5427 \\
\bottomrule
\end{tabular}
\end{table}

\noindent
\textsc{full} exceeds \textsc{cond} by 0.085, approximately 2.3 pooled standard
deviations, with ranges that barely overlap. Read alone, this appears to show
LOB features contributing.

It does not, and the distinction matters. \textsc{cond} is a model trained with its dominant input stream zeroed, a
damaged network rather than a fair alternative predictor. The lookup table comparison in Table~\ref{tab:baseline} is the
correct counterfactual, and it shows no gain. The ablation establishes that the
model's internal representation depends on the LOB window; it does not establish
that the window carries information unavailable elsewhere. We flag this because
the ablation was run first, and taken alone would have supported the opposite
conclusion.

\textsc{lob} at $0.5034$ confirms the defect reproduces to chance across five
seeds.

\subsection{Status of the reinforcement learning component}
\label{sec:rl-status}

Version~1 proposed a latency-aware PPO agent consuming fill probability
estimates, with a reward function penalising adverse fills attributable to
execution delay. We report no RL results in this version, for two reasons.

The agent was trained only on synthetic data, where it did not converge:
per-episode PnL oscillated without trend across 2{,}000 episodes. It also consumed fill probability estimates from the defective model, which
makes those runs uninterpretable regardless of the environment.

The reward function itself remains, in our view, a reasonable contribution, but
it is untested, and we now present it as a proposal rather than a validated
component. Given the results above, an agent conditioned on our fill probability
estimates would receive no information beyond quote offset and latency, both of
which it already selects directly.

% ══════════════════════════════════════════════════════════════════════════════
\section{Interpretation}
\label{sec:interpretation}

Three explanations are consistent with our results. We cannot fully distinguish
them, and say so.

\subsection{Sampling frequency (most likely)}

At a median 18.2-second inter-update interval, a 200\,ms latency window is
roughly $1/90$th of the typical gap between consecutive observations. Nearly all short-horizon dynamics (queue position changes, fleeting imbalance,
order flow bursts) occur entirely between our samples and are unobservable.

This is consistent with the wider literature. \citet{arroyo2024} and
\citet{maglaras2022} demonstrate meaningful fill prediction from LOB data at
tick resolution, three to four orders of magnitude finer. Our null does not contradict those findings. It locates a frequency below which
the signal they exploit is no longer present.

If correct, this has a direct practical implication that inverts the framing of
Version~1. The binding constraint for a retail participant is data frequency rather than
model sophistication or dataset size. And data frequency is purchasable:
full-depth or tick-level feeds are commercially available. The
democratisation argument weakens accordingly, since the barrier turns out to be
another form of infrastructure cost rather than something a model can substitute
for.

\subsection{Feature construction}

Our 109 features are standard but were computed under parameters adapted for
sparse data (Section~\ref{sec:data}). A rolling window of 10 snapshots spans
roughly seven minutes of wall-clock time here. It is possible that better-designed features (explicitly time-aware rather than
count-based, or incorporating trade prints alongside book snapshots) would
recover signal this construction destroys. We consider this less likely than the frequency
explanation, since it would have to overcome a gap of $0.0003$, but we have not
excluded it.

\subsection{Architecture and objective}

The survival formulation with 90{,}754 parameters may be poorly matched to this
data. The ETH result, where the model actively underperforms a table, shows the
architecture can destroy information at small sample sizes. A simpler
approach, gradient boosting on flat features being the obvious candidate, might
extract more. We regard this as a genuine gap in the present work: we did not
run that comparison, and it is the first thing a referee should ask for.

% ══════════════════════════════════════════════════════════════════════════════
\section{Limitations}

We state the limitations of the negative result directly, since a null claim is
only as strong as the effort made to overturn it.

\textbf{Single split.} All results use one 70/15/15 temporal split. Repeated
seeds capture initialisation and batch-order variance but not sampling
variance. A different split might yield different numbers, though the
$+0.0003$ pooled gap would need to move by two orders of magnitude to change the
conclusion.

\textbf{One architecture.} We tested an LSTM survival model. The
convolutional-transformer of Version~1 was not re-evaluated after the
conditioning fix, and no gradient-boosted or linear baseline was run. We can say
this architecture does not beat a lookup table; we cannot say no model would.

\textbf{One venue, one feed type.} Kraken's depth-10 book channel. Neither
other exchanges nor other feed types (trade prints, full-depth, message-level)
were tested, and the frequency hypothesis in
Section~\ref{sec:interpretation} predicts precisely that a different feed would
give a different answer.

\textbf{Simulated fills.} Labels come from replaying subsequent book states, not
from executed orders. Queue position is not modelled, so we implicitly assume
a quote at a price level fills when that level trades through, which is
optimistic for a latecomer to the queue. Whether this biases against microstructure features
specifically is unclear.

\textbf{ETH undertrained.} We report the ETH result as-is rather than tuning
it. It should be read as a statement about this architecture at this sample
size, not about ETH/USD.

\textbf{No live trading.} Everything here is offline evaluation. No capital was
deployed and no strategy is validated for deployment.

% ══════════════════════════════════════════════════════════════════════════════
\section{Implications}

\subsection{For researchers}

Report the trivial baseline. Our ablation showed a clean $+0.085$ separation and
would, presented alone, have supported the opposite conclusion of this paper.
The lookup table cost two minutes to implement and reversed the finding. In microstructure prediction, where a small number of variables often carry most
of the mechanically available signal, an ablation against a damaged version of
one's own model is not sufficient evidence that features matter.

Sampling frequency deserves explicit treatment. Papers reporting LOB prediction
results do not consistently state inter-update intervals, yet our results
suggest this parameter can determine whether an effect exists at all.

\subsection{For retail practitioners}

The direct implication is favourable in one narrow sense: if fill probability at
these frequencies is a function of quote offset, latency, and side, it can be
estimated with a lookup table from a few weeks of data. No GPU, no deep
learning, no ML expertise. The model was not necessary.

The broader implication is unfavourable. Version~1 argued that AI could
substitute for infrastructure and thereby democratise liquidity provision. Our
results indicate that at the data frequency available for free, there is no
microstructure edge to extract. Accessing finer data means paying for it, which
reintroduces the infrastructure barrier the framework was meant to circumvent.

\subsection{For market design}

Our finding is orthogonal to the case for batch auctions and speed bumps
\citep{budish2015,aquilina2022}, which rest on the social cost of the speed
race rather than on whether individual participants can compensate
technologically. If anything it strengthens the case for venue-level reform: we
find no evidence that participant-level technical adaptation is a viable
substitute at retail data access, which removes one argument against structural
intervention.

We are more cautious than Version~1 regarding AI-driven democratisation of
market microstructure. That version asserted the structural concentration of
liquidity provision is ``a solvable problem, and artificial intelligence is part
of the solution.'' We have no evidence for that claim and some against it.

% ══════════════════════════════════════════════════════════════════════════════
\section{Conclusion}

We set out to test whether machine learning models trained on limit order book
microstructure could substitute for execution speed in market making, enabling
participation by traders operating at 200\,ms--2\,s latencies. Using 269{,}330
LOB snapshots collected over 70 days across three cryptocurrency pairs, we find
they do not, at the data frequency available from a public exchange feed.

A survival-analysis model over 109 microstructure features matches a 32-cell
lookup table on quote offset, latency, and side to within 0.0003 AUC pooled, and
fails to beat it on any of three assets individually. A controlled ablation
places essentially all measurable predictive content in quote geometry. A
deliberately reproduced conditioning defect performs at chance across five
seeds, confirming the corrected pipeline behaves as intended.

We regard the most probable explanation as sampling frequency: at a median
18.2-second inter-update interval, the short-horizon dynamics that
higher-frequency work exploits occur between our observations. If so, the binding constraint for retail participation is data access rather
than modelling capability. That is a less encouraging conclusion than the one we
set out to reach, and a more actionable one.

The framework of Section~\ref{sec:formal} remains, we believe, a correct
formalisation of the problem. What we have shown is that its first viability
condition fails in the specific setting where it would matter most to the
participants it was designed to serve. Establishing whether it holds at finer
data resolution is the natural next step, and requires only a feed we did not
have.

% ══════════════════════════════════════════════════════════════════════════════
\section*{Data and Code Availability}
\addcontentsline{toc}{section}{Data and Code Availability}

The complete collection pipeline, feature engineering code, models, ablation and
baseline scripts, and result files are available at
\url{https://github.com/azure-cyber/latency-compensating-mm}. All results in
this paper are reproducible from that repository. The diagnostic and baseline
scripts (\texttt{diagnose.py}, \texttt{baseline\_check.py},
\texttt{ablation.py}, \texttt{per\_asset.py}) are included so the null result can
be checked directly.

\section*{Disclosure of AI Assistance}
\addcontentsline{toc}{section}{Disclosure of AI Assistance}

The author used Claude (Anthropic) extensively throughout this project: for
implementation of the data collection pipeline, feature engineering, models, and
evaluation code; for diagnostic analysis including identification of the
conditioning defect described in Section~\ref{sec:defect}; and for drafting and
revision of this manuscript. The author designed the study, directed the
analysis, made all substantive research decisions, verified the results, and is
solely responsible for the content and any errors.

\section*{Acknowledgements}
\addcontentsline{toc}{section}{Acknowledgements}

This work was conducted independently, without institutional affiliation,
funding, or supervision. The author declares no conflicts of interest and holds
no positions in any asset discussed.

% ══════════════════════════════════════════════════════════════════════════════
\bibliography{refs}

\end{document}
